\section{Introduction and Context}
\label{sec:introduction}

Intrusion detection is an operational problem as well as a classification problem.
Missed attacks can leave harmful activity unchallenged, while false alarms consume
analyst attention and make genuine incidents harder to prioritise. Practical guidance
therefore treats alert validation, false-positive identification, and threat
prioritisation as central parts of intrusion monitoring~\citep{scarfone2007idps}.

Networks produce a continuous stream of connections whose meaning is rarely
contained in a single record. A destination port, a packet count, or a flow duration can
be informative, but the surrounding communication pattern often matters just as much.
A distributed denial-of-service attack may appear as many sources converging on one
destination, whereas a scan may appear as one source contacting many hosts or services.
This relational character makes network traffic a natural candidate for graph-based
analysis. Network events also carry semantic cues: protocol and port
combinations, traffic magnitudes, and recurring behavioural descriptions can suggest
what an interaction represents. An intrusion detection system that sees only one of
these views is working with an incomplete account of the traffic.

This report investigates whether structural and semantic evidence can be combined in a
single edge-level intrusion detection framework using NF-UNSW-NB15 and NF-ToN-IoT. One
experimental boundary must be stated from the outset: the current graph construction
aggregates flows by source IP, destination IP, and attack class. Ground-truth labels
therefore participate in defining edge identity, even though the attack name is excluded
from the model's input features and generated text. The
resulting scores compare representations and interaction mechanisms under a controlled
protocol, but they are not deployment-valid estimates of detection performance on unseen,
unlabelled traffic.

NF-UNSW-NB15 derives from UNSW-NB15, whose original benchmark was created from a mixture
of normal activity and contemporary attacks in the UNSW Canberra Cyber Range
Laboratory~\citep{moustafa2015unsw}. NF-ToN-IoT derives from the TON\_IoT datasets and
provides an IoT-focused telemetry setting~\citep{moustafa2021toniot}. This study uses both
in their NetFlow-standardised forms~\citep{sarhan2021netflow}. One identical processing
and modelling pipeline is applied to the two datasets. Every architectural choice is
shared between them. Only two quantities selected on validation data differ.

The two datasets were chosen because they occupy very different points on the
class-imbalance axis. After aggregation, NF-UNSW-NB15 contains 656 edges, is 47.4\%
benign, and has a most-to-least-frequent class ratio of $12.4{:}1$. Eight of its nine
attack classes contain 40 edges and the ninth contains 25, so just one falls at or below
that threshold. NF-ToN-IoT contains 2{,}127 edges, is 82.1\% benign, and has a most-to-least-frequent
class ratio of $582{:}1$. Seven of its nine attack classes contain 35 edges or fewer. Its
distribution is far more skewed. By the most-to-least class ratio, the two
graphs differ by roughly a factor of 47 in imbalance severity.

That difference is the methodological reason for using two datasets. A single dataset
cannot separate a property of the architecture from a property of the class distribution
that architecture happened to face. Holding the architecture fixed while changing the
imbalance regime by roughly $47\times$ makes that distinction experimentally addressable.
RQ4 already asks how class imbalance affects fusion and feedback, and the two-dataset design
gives that question two substantially different regimes to compare. No stage is assumed in
advance to be sensitive or insensitive to that change.

The scoring protocol also differs in one necessary respect. Macro-F1 covers ten classes on
NF-UNSW-NB15 and eight classes on NF-ToN-IoT, because two NF-ToN-IoT classes are too small
to survive five-fold splitting. Their edges remain in the graph for message passing but
are excluded from the metric. Because the class counts differ, absolute scores are not
comparable across the two datasets. Only the shape of the model ladder is. This constraint
applies throughout the report.

\subsection{Why combine graph and language representations?}

Graph neural networks (GNNs) provide a direct way to represent communication structure.
In this work, IP addresses form the nodes of a directed graph and aggregated flows form
labelled edges. A graph-attention encoder propagates information between communicating
hosts and produces a structural representation for each edge. This branch can learn
patterns that are difficult to express as independent rows in a conventional feature
table. It does not, however, possess an explicit vocabulary for interpreting the traffic.
The graph tells the model \emph{who communicated with whom} and how that interaction is
positioned in the network. It does not explain what a high-frequency TCP
connection to a particular service may signify.

A pretrained cybersecurity language encoder offers a complementary view. The same
aggregated flow can be written as a short, controlled sentence describing its endpoints,
protocol, port, frequency, volume, and duration. The implementation uses CySecBERT, a
BERT-class model adapted to cybersecurity text and evaluated by its authors on
domain-specific representation and classification tasks~\citep{bayer2024cysecbert}. It is
frozen and used purely as a sentence encoder. Throughout this report, this component is
called the \emph{semantic branch}. That terminology distinguishes its representational
role from that of a generative large language model.

Neither representation is assumed to win in every case; the design instead lets them
challenge and complement one another. The parallel fusion stage runs
the graph and semantic encoders side by side and combines their edge embeddings through
Adaptive Gated Attention Fusion, referred to throughout as the fusion model. The feedback stage adds uncertainty-guided
semantic feedback: edges with high-entropy GNN predictions are selected for semantic
consultation, and the resulting signal is reintroduced by modifying the edge
representations that the next graph update consumes. A second design, which biases graph
attention directly, was also implemented and measured. It produced no change in prediction
at all, and is reported as a verified negative result.

The coupling is bidirectional in routing. The graph model's per-edge uncertainty decides
which edges are consulted, while the semantic branch returns a per-edge judgement that
re-enters the edge representations consumed by the next graph update. It is not
bidirectional in content: the consultant's judgement depends on that edge's text alone,
not on the graph model's current belief about the edge.

This general dual-stream idea is not novel by itself. Parallel BERT and graph encoders
with attention-based fusion have already been studied outside network security, for
example in misinformation detection~\citep{moorthy2025dualstream}. Work in network
intrusion detection has also combined GNNs and generative language models for a different
role: XG-NID performs intrusion classification with a heterogeneous GNN and uses an LLM
after classification to generate explanations and remedial guidance
\citep{farrukh2025xgnid}. DAS is the closest prior work to the feedback mechanism
considered here. It already implements a closed feedback loop between a frozen GNN and an
LLM, feeds refined semantics back to update the same graph learner, and constructs node
descriptions by expressing structural statistics in natural language. Its setting is
node-level and concerns general text-rich graphs~\citep{thapaliya2025das}. The
contribution here is therefore instantiation plus measurement: an edge-level
network-intrusion-detection instantiation, together with a diagnostic question about
whether a constrained semantic consultant exploits a feedback channel that has the
capacity to influence the next graph update. Within the literature reviewed here, we are
not aware of prior work combining bidirectional iterative coupling between a graph model
and a semantic branch, joint attention-based fusion for the final classification, and
edge-level network intrusion detection.

\subsection{Research problem and questions}

The research problem is not simply whether adding a language embedding raises a score.
If one branch is already a strong classifier, a fusion module can appear successful while
doing little more than copying that branch. Likewise, a feedback loop can improve the
final output because semantic features are fused at the classifier, even when the claimed
attention re-entry mechanism contributes nothing. The evaluation must separate
the value of the two modalities, the fusion mechanism, the uncertainty-selection policy,
and the iterative update itself.

Five research questions frame the study:

\begin{description}
    \item[RQ1.] How do the structural GNN branch and the semantic branch perform
    individually on NF-UNSW-NB15 and NF-ToN-IoT under the same five-fold out-of-fold
    evaluation protocol?

    \item[RQ2.] Does the fusion model score above both single-branch rungs on NF-UNSW-NB15 and
    NF-ToN-IoT, and are those differences separated under the bootstrap protocol?

    \item[RQ3.] Does uncertainty-guided semantic feedback add value beyond static fusion
    on NF-UNSW-NB15 and NF-ToN-IoT? If so, does the gain come from semantic re-entry into
    the graph representation or from final-output fusion?

    \item[RQ4.] How do graph topology, repeated endpoint pairs, class imbalance, and the
    strength of the semantic consultant affect fusion and feedback on NF-UNSW-NB15 and
    NF-ToN-IoT?

    \item[RQ5.] How does the feedback model compare against published graph-based
    intrusion detection architectures re-trained on the same aggregated representation
    under the same protocol on NF-UNSW-NB15 and NF-ToN-IoT, and how much of any observed
    difference is attributable to representation and tuning rather than to architecture?
\end{description}

These questions deliberately allow a negative answer. The experiments are meant to determine
when the proposed interaction works, when it does not, and why. A model ladder of GNN, semantic
encoder, fusion and then feedback is only meaningful if each higher rung is compared fairly
with the strongest relevant lower rung.

\subsection{Objectives and contributions}

The practical objective is to build one reproducible pipeline that transforms
NF-UNSW-NB15 and NF-ToN-IoT NetFlow records into two aligned views and evaluates every
stage under one consistent protocol. The work makes five contributions.

\begin{enumerate}
    \item It implements a directed, edge-classification graph pipeline in which flows are
    aggregated by source IP, destination IP, and attack class, with centrality-based node
    features and statistical edge attributes.

    \item It constructs a corresponding internal knowledge representation and serialises
    each edge as controlled natural language. The text supplied to the encoder is kept
    label-free by an explicit validation guard, preventing the semantic branch from seeing
    the target attack name or a synonymous relation label.

    \item It implements and compares two multimodal stages: the fusion model for parallel structural--
    semantic fusion, and uncertainty-guided semantic feedback that reintroduces semantic
    evidence into the edge representations of uncertain edges. The feedback model consults
    a classification head trained on the semantic embeddings. The feedback model is reported under that
    consultant and under the weaker embedding-only scorer it was first built with, so the
    reader can see how much of the result is the mechanism and how much is the consultant.

    \item It provides a controlled evaluation on NF-UNSW-NB15 and NF-ToN-IoT, including
    unimodal baselines, fusion comparisons, uncertainty-selection ablations, trained-head
    baselines, mechanism-isolation tests, and two published graph-based intrusion
    detection architectures re-trained on the same representation. It reports both
    positive and negative results to show which components carry the measured gain,
    including five measured attempts to make an embedding-only consultant's advice usable
    by the feedback channel, none of which succeeded.

    \item It provides a cross-dataset comparison under one shared architecture across two
    substantially different class-imbalance regimes, testing whether the behaviour of each
    stage in the model ladder is a property of the method or of the class distribution it
    faces.
\end{enumerate}

The contribution is not presented as a universal improvement over existing
intrusion detectors. It is an implementation and controlled evaluation of a particular
form of structural--semantic cooperation, together with evidence about its limits. The
headline of that evidence is stated here so it is not mistaken for a footnote later. With a
trained semantic consultant the feedback model is the highest rung on both datasets, separated
from the GNN model, from the fusion model and from every faithful re-trained baseline. With
the weaker consultant it was first built with, none of those comparisons is separated at all.
And when final-output fusion is removed, so that semantic evidence can reach the classifier
only by re-entering the graph, neither consultant produces a gain we could separate from zero.
The channel is live and measurable, and what the re-entry itself adds is, on these graphs,
nothing we could distinguish from nothing.
Deployment-valid evaluation remains open and requires graph construction in which edge
identity can be determined without access to the target class.

\subsection{Report structure}

The report is organised into five sections. This Introduction defines the problem,
experimental scope, research questions, and contribution. Section~\ref{sec:related}
reviews graph-based intrusion detection, cybersecurity language encoders,
structural--semantic fusion, and iterative graph--language feedback, and positions the
study against the closest prior mechanisms. Section~\ref{sec:method} describes the two
datasets, graph and text construction, the structural and semantic branches, the fusion model,
uncertainty-guided semantic feedback, the baseline protocol, and the shared evaluation
procedure.

The Results section reports the predefined model ladder, baseline comparisons, ablations,
and mechanism-isolation measurements for both NF-UNSW-NB15 and NF-ToN-IoT, while keeping
the different scored class sets explicit. The Discussion section interprets those
measurements against the research questions, including the cross-dataset imbalance
comparison and the distinction between a feedback channel having capacity and a
constrained semantic consultant making use of that channel. It also returns to the
controlled-protocol limitation and the work still required for deployment-valid
intrusion detection.
